\section{Background}
\label{sec:background}

% Here we detail the previous work done in this field and how we build on it.
Spectre attacks exploit the discrepancy between architectural correctness and 
micro-architectural state, allowing transiently executed instructions to leak 
sensitive information through cache side channels \cite{SpectreBase}. Existing 
defenses have been evaluated extensively using architectural simulators such 
as gem5~\cite{Narayan:HFI:2024}, but simulation abstractions may not fully 
capture the pipeline timing, cache replacement behavior, and hardware resource 
contention of synthesizable RISC-V processors. In evaluating different software 
defenses in SpectreShield~\cite{Khan:SpectreDef:2025}, gem5 was used as the main 
simulation tool to observe speculation behavior on micro-architectural elements. 
Here, they explicitly list out the various metrics used to evaluate Spectre in 
cache contention, branch prediction accuracy, DRAM access \& power variability, 
and secret key recovery rates.

Prior research has implemented speculative-execution defenses in both gem5 and 
RISC-V RTL for caches (SpecLFB) \cite{Cheng:SpecLFB:2024} and taint-tracking 
(Secure Speculative Execution) \cite{sabbagh2021secure}. In addition, SecureBOOM 
semi-automates integrating security features into the microarchitecture without 
having the designer consider every potential security gap; this is done through 
formal analysis on the RTL design, which was implemented on the Berkeley 
Out-of-Order Machine (BOOM). Important questions remain regarding the consistency 
of their security guarantees and performance overhead across these abstraction 
levels \cite{Jauch:SecureBOOM:2023}. Indeed, $\mu$STT mentions how SecureBOOM's 
infrastructure negatively impacts area overhead \cite{Chen:uSTT:2025}.

This project implements an existing speculation-aware defense, 
$\mu$STT~\cite{Chen:uSTT:2025}, in an open-source RISC-V processor, and compares 
our results with those published for its gem5 and standalone RTL evaluations. 
By measuring security, performance, and hardware cost in a complete core, we aim 
to determine whether the published trends hold in a full RTL design, and to 
characterize the trade-off between security and processor efficiency.
